% Literature search: insecure-code.
% Author's note (verification and reproduction), not compiled text:
%   scholar session "insecure-code" (searched 2026-11-27, 09:51-14:24 UTC);
%   full query strings and the query x provider x hits matrix in
%   litteratursokning/insecure-code-queries.md (scripts
%   insecure-code-queries-ph.sh, -redo.sh, -repair.sh); progress and
%   screening log in litteratursokning/insecure-code-progress.md; screen
%   sheets, rest sample (seed 20261127) and counter sample (seed 20261128)
%   in litteratursokning/insecure-code-screen/; check by meaning in
%   litteratursokning/insecure-code-check/; bibliography with provenance
%   blocks in litteratursokning/insecure-code.bib; raw export
%   litteratursokning/insecure-code-export.csv.
%   Screened by an LLM classifier and an AI agent; no decision has been
%   reviewed by the author, so the hit list shows no row as cited yet.

\chapter{Do people who program with an LLM-based coding assistant produce
  less secure code, and are they more confident in its security, than
  people without one?}
\label{insecure-code:chapter}

\emph{Status: screened by an AI agent; awaiting human review.}

The notes warn that code written with an AI assistant may be less secure
than it looks.
The claim is most often repeated in the form of one experiment, in which
participants with an AI assistant wrote significantly less secure code than
those without one and were more likely to believe that their code was
secure \autocite{Perry2023Users}.
This chapter asks whether that holds as a general statement, in two parts
kept apart.
The first concerns people: is code written by people with an assistant less
secure than code written by people without one, and are they more confident
in it?
The second concerns the tools: how often does code generated by language
models, coding assistants and coding agents contain vulnerabilities, and has
that changed with newer models and with agents?
A third, short strand asks how often models suggest software packages that
do not exist, and whether attackers have exploited that.
The strength tested is a comparison (\enquote{less secure than without}) and
a frequency (\enquote{often}), for empirical work from 2021 to 2026.

\section{Method}
\label{insecure-code:method}

\paragraph{Databases and their health.}
All searches were run on 27~November 2026 through the databases' programming
interfaces, without an upper limit on the number of results, in
OpenAlex, Web of Science, Scopus, IEEE Xplore, Semantic Scholar, arXiv and
DBLP.
DBLP answered every query with a page that stops automated clients, so it
could not be searched.
The university network route was unavailable for the first searches, which
Scopus therefore answered in its standard view without abstracts; it came
back in the afternoon, and the repeated and repair searches went through it.
IEEE Xplore had spent its daily quota before 9 of the 18 queries, Semantic
Scholar refused 4 as rate-limited even when repeated, and arXiv answered 8
with server errors on long queries; its answer of 0 to C5 is doubtful,
since a subset of C5 (S4) returned 882 records, and is treated as not
searched.
These cells are marked as failed in \cref{insecure-code:tab:queries}: they
were not searched and are not zeros.
Semantic Scholar answered through its bulk search.

\paragraph{Queries.}
Each concept was searched in several wordings, fixed before the queries were
run: the assistant (\enquote{large language model}, LLM, \enquote{language
model}, \enquote{AI assistant}, \enquote{code assistant}, \enquote{coding
assistant}, \enquote{programming assistant}, \enquote{AI pair programmer},
Copilot, ChatGPT, Codex, \enquote{generative AI}, \enquote{AI-generated},
\enquote{AI-assisted}, \enquote{coding agent}, \enquote{vibe coding},
singular and plural); security (security, secure, insecure, vulnerability,
vulnerable, CWE, \enquote{security weakness}, \enquote{security flaws},
\enquote{code security}); the kind of study (\enquote{user study},
participants, \enquote{controlled experiment}, \enquote{between-subjects},
\enquote{lab study}); confidence (confidence, overconfidence, trust,
overtrust, overreliance, perception); generated code (\enquote{generated
code}, \enquote{code generation}, \enquote{code suggestions}, \enquote{code
completion}, and in S9 the verb and noun forms generate, generating,
generation, write, written, with the objects code, applications, programs
and benchmarks); code in the wild (GitHub, \enquote{open source},
repositories, commits, \enquote{pull requests}); the population (students,
novices, beginners, \enquote{introductory programming}, CS1); agents
(\enquote{coding agent}, \enquote{AI agent}, \enquote{LLM agent}, agentic,
\enquote{reasoning model}); and package hallucination (\enquote{package
hallucination}, \enquote{hallucinated package}, \enquote{non-existent
package}, \enquote{phantom dependencies}, slopsquatting, and the described
form: hallucination near package, dependency or library near npm, PyPI,
\enquote{supply chain} or typosquatting).
Each counter-query keeps every concept term of its supporting query and adds
the terms of the opposite finding (\enquote{no difference}, \enquote{no
significant}, \enquote{more secure}, \enquote{fewer vulnerabilities}, safer,
improve, replication, comparable, \enquote{human-written}, newer,
longitudinal), or terms of distrust and calibration (C2), of repairing
vulnerabilities (C4, C5), or of low, falling, mitigated or unexploited
package hallucination (C8).
\Cref{insecure-code:tab:queries} gives the queries and the number of records
each database returned.

% Cells copied from the session's search records (scholar sessions show
% insecure-code); full query strings in insecure-code-queries.md.
\begin{table}[tbp]
  \begin{sidecaption}{Queries and hits per database, 27~November 2026.
    S = supporting query, C = counter-query.
    Databases: OpenAlex (OA; not open access), Web of Science (WoS),
    Scopus, IEEE Xplore (IEEE), Semantic Scholar (S2), arXiv.
    The searches have no upper limit: the numbers are all hits.
    \enquote{--} means that the search failed: IEEE Xplore's daily quota was
    spent, Semantic Scholar refused the request as rate-limited, or arXiv
    answered with a server error.
    \enquote{0*}: an answer of 0 that is treated as not searched (see the
    text).
    DBLP failed on every query and could not be searched; it has no column.
    WoS received the keywords as \texttt{TS=(\dots)}, Scopus as
    \texttt{TITLE-ABS-KEY(\dots)}, the others as Boolean free text, in title
    and abstract where the database allows it.}[insecure-code:tab:queries]
    \centering\footnotesize
    \begin{tabular}{@{}lp{0.36\linewidth}rrrrrr@{}}
      \toprule
      & Concepts & OA & WoS & Scopus & IEEE & S2 & arXiv \\
      \midrule
      S1 & assistant, security, user study & 1884 & 510 & 792 & 433 & -- & 463 \\
      S2 & assistant, security, code, confidence or trust & 2134 & 343 & 537 & 835 & -- & 728 \\
      S3 & assistant, vulnerability, generated code & 1187 & 340 & 541 & 417 & -- & 447 \\
      S4 & agent, vulnerability, code & 3572 & 120 & 252 & 194 & 2338 & 882 \\
      S5 & assistant, vulnerability, code in the wild & 3498 & 960 & 1744 & 1022 & 2216 & -- \\
      S6 & assistant, security, code, students or novices & 1303 & 257 & 418 & 742 & 582 & 190 \\
      S7 & package hallucination (named) & 97 & 13 & 23 & -- & 53 & 30 \\
      S8 & assistant, hallucination, package, registry or supply chain & 80 & 11 & 21 & -- & 42 & 24 \\
      S9 & assistant, vulnerability, generate or write, code or applications or benchmarks & 5630 & 1535 & 2570 & -- & -- & -- \\
      \midrule
      C1 & S1 and counter terms & 952 & 266 & 284 & 197 & 786 & -- \\
      C2 & S2 and distrust or calibration & 678 & 70 & 110 & 204 & 254 & 226 \\
      C3 & S3 and counter terms & 647 & 184 & 202 & 238 & 701 & -- \\
      C4 & assistant, vulnerability, repair or fix & 4972 & 1391 & 1881 & -- & 3326 & -- \\
      C5 & S4 and repair or counter terms & 2467 & 81 & 135 & -- & 1854 & 0* \\
      C6 & S5 and counter terms & 1800 & 504 & 580 & -- & 1613 & -- \\
      C7 & S6 and counter terms & 762 & 136 & 115 & -- & 449 & 164 \\
      C8 & S7 or S8, and rare, declined, mitigated, newer, exploited or registered & 929 & 221 & 247 & -- & 583 & -- \\
      C9 & S9 and counter terms & 3016 & 819 & 897 & -- & 2814 & -- \\
      \bottomrule
    \end{tabular}
  \end{sidecaption}
\end{table}

\paragraph{Studies the queries had to return.}
Four studies were named before the screen as ones the queries must return:
the two early user studies \autocite{Perry2023Users,Sandoval2023LostC}, the
first measurement of Copilot's code \autocite{Pearce2022Asleep} and the
first large measurement of package hallucination
\autocite{Spracklen2025Package}.
All four were returned, each by several queries and databases.

\paragraph{Check by meaning.}
After the keyword searches, six questions in plain words were put to
OpenAlex's semantic search (at most 50 works each), and every returned work
was looked up in the session: whether people with an assistant write less
secure code and are overconfident; whether studies found no difference; how
often generated code is vulnerable; whether newer models or agents generate
more secure code; vulnerabilities introduced by coding agents in real
repositories; and package hallucination.
For the package-hallucination question 35 of the 50 works were in no
session, because the queries then had no such concept; S7, S8 and C8 were
added and run on every database.
For the others, 3 to 19 works were in no session, most of them overviews,
theses in other languages, or studies of agents' pull requests that measure
no security outcome; but six bore on the second sub-question, among them a
secure-coding benchmark \autocite{Bhatt2023CyberSecEval} and studies of
generated web applications.
Their abstracts say that models \enquote{generate} insecure code or
applications, or speak of benchmarks, rather than of \enquote{generated
code} or \enquote{code generation}; the generated-code concept was
therefore searched again in its verb and noun forms (S9, C9), after which
three of the six were in the session.
The other three describe security in words that the vulnerability block
of the generated-code queries does not hold (\enquote{security best
practice}, \enquote{secure authentication}, \enquote{authorization
exposure}) and remain outside it.
A ranked list of this kind has no count and no stopping point, so it was
used only to check the keyword queries.

\paragraph{Screening.}
The 18 queries returned 21,436 distinct records.
The first tier, screened in full, was every record at place 100 or better
in a list that a database ordered by relevance: 2,644 records.
The 3,008 records placed only in answers that are not ordered by relevance
(those of Web of Science and Scopus) fall outside it and were reached only
through the samples.
A language model (Claude Sonnet~5.5, without examples, against a written
statement of the question and of the exclusion criteria below) classified
each record by title and abstract, after missing abstracts had been filled
in from other databases where they held one.
The AI agent then read every record the model marked as qualifying, the
titles of the records it kept with low confidence, and the full texts of
every source cited below; it overturned four of the model's keeps and
recorded each cited source with its theme.
Of the 2,644, 258 were kept as supporting and 66 as qualifying (12.3\,\%
bearing).
Of the 20,595 records outside the first tier, two seeded random samples of
60 were screened like it: one from all of them (seed 20261127), screened by
the model, held 2 bearing records (3.3\,\%; 95\,\% modified Wilson interval
0.6 to 11.4\,\%); one from the 14,126 returned by a counter-query (seed
20261128), screened by the agent, held 1 (1.7\,\%; 0.1 to 8.9\,\%), a
comparison of generated and human-written smart contracts whose abstract
does not state the direction of the result.
The rest were not screened.
No decision has been reviewed by the author, and no seeded sample of the
machine's decisions has yet been reviewed, so the machine's error rate is
unknown.

\paragraph{Exclusion criteria.}
A record is excluded if it is one of the following.
\begin{description}
  \item[No abstract] It has no abstract, and none could be added from its
    DOI or other identifier.
  \item[Detection only] It uses language models to detect, explain or
    repair vulnerabilities in existing code and measures nothing about the
    security of code that models or assistants generate.
  \item[New technique] It proposes a new technique for generating or
    repairing code securely and reports mainly its own improvement, not a
    measurement.
  \item[Attack on model] It attacks the model or agent itself (prompt
    injection, jailbreak, poisoning) without measuring the security of the
    code from ordinary use.
  \item[No security outcome] It studies AI coding tools but measures no
    security outcome.
  \item[No evidence] It gives an opinion, overview or position without
    evidence of its own or a stated method.
  \item[Off topic] It belongs to another field and shares keywords only.
\end{description}
Among the 2,432 excluded records, the criteria were given 910 times as off
topic, 477 as detection only, 327 as attack on model, 314 as new technique,
263 as no evidence, 150 as no security outcome and 26 as no abstract (a
record can carry two).
After the screen, missing abstracts were sought again in Scopus through the
university route; 187 records gained one, 27 of them records already
excluded for lack of an abstract, which were re-read and stayed excluded
under another criterion.

\paragraph{Verification and deduplication.}
Every source cited below was opened.
Full texts were read for the open-access ones (preprint servers, the
authors' or a university repository) and, once the university route was
back, for one publisher's chapter; seven closed-access papers, behind
publishers that refuse automated access, could be read only in their
abstracts and are marked as such where they are cited.
The search software merges copies of a work within a search (same DOI, else the same
title and authors) and folds a preprint into its published version when the
match is unambiguous; copies that reached the session from different
searches were kept as separate records and given the same decision, as
PRISMA-S asks a report to describe \autocite{Rethlefsen2021PRISMAS}.

\section{Results}
\label{insecure-code:results}

\paragraph{People with and without an assistant.}
The controlled comparisons disagree.
In the experiment the claim comes from, 47 participants, 33 with and 14
without an assistant based on OpenAI's codex-davinci-002 model, solved five
security-related tasks; those \enquote{who had access to an AI assistant
wrote significantly less secure code than those without access to an
assistant} \autocite[abstract]{Perry2023Users}.
The authors note that the participants \enquote{consisted mainly of
university students} \autocite[sect.~7.2]{Perry2023Users}.
A randomised study of 58 computer-science students writing a linked list in
C with an older Codex model found that \enquote{AI-assisted users produce
critical security bugs at a rate no greater than 10\% more than the
control} \autocite[abstract]{Sandoval2023LostC}.
In a replication with Copilot and 25 participants, 19 of them graduate
students, \enquote{access to Copilot accompanies a more secure solution when
tackling harder problems}, with no effect on the easier problem
\autocite[abstract]{Asare2024UserCentered}.
In a within-subject study of 7 novices and 6 professionals, Copilot
\enquote{did not have an effect on code security between novices and
experienced developers}, though it reduced certain kinds of vulnerability
\autocite[abstract]{ZareShahneh2025Copilot}; the peer-reviewed version adds
that most participants wrote insecure code regardless of experience
\autocite{ZareAssal2026Copilot}.
The largest comparison found, 159 freelance developers assigned to no AI or
to the free or the paid version of Google's Gemini, observed no significant
differences in code security between the conditions, while programming
experience improved security \autocite[abstract]{Jost2026Gemini}.
A password-storage study with 21 students and 80 freelancers found that many
relied on AI-generated code and prioritised convenience, and that
\enquote{AI tools alone are insufficient without targeted guidance}
\autocite{Yardim2026Password} (abstract only).
Among computing students, AI assistance had \enquote{a high chance of
producing insecure programs} in common curriculum tasks
\autocite{Yuan2025Mindset} (abstract only), and 61 undergraduates using a
deliberately poisoned model \enquote{frequently accepted insecure
LLM-generated code despite prior security coursework and stated skepticism}
\autocite[abstract]{Park2025Aware}.

\paragraph{Confidence.}
Over-confidence has been measured directly in one comparison with a control
group: participants with an assistant \enquote{were also more likely to
believe they wrote secure code} \autocite[abstract]{Perry2023Users}.
A later observational study of 100 participants, all with AI suggestions,
found that they could not tell the most from the least secure suggestion,
yet their \enquote{confidence in their code's security correlated with its
actual security} \autocite[abstract]{Khalid2026Trust}.
With an assistant manipulated to suggest insecure password storage,
freelance developers showed \enquote{a high level of trust in AI-generated
code, even when insecure suggestions were presented}
\autocite{Serafini2025Manipulated} (abstract only).
In a survey of 105 developers, \enquote{only about a quarter of the
developers expressed confidence in the code generated by AI}
\autocite[abstract]{Kudriavtseva2025Perceptions}.

\paragraph{How often generated code is vulnerable.}
Of 1,689 programs that Copilot generated in 2021 for 89 scenarios built
around dangerous weaknesses, \enquote{approximately 40\,\% [were]
vulnerable} \autocite[abstract]{Pearce2022Asleep}.
A replication with a newer Copilot found that \enquote{the percentage of
vulnerable code suggestions has reduced from 36.54\% to 27.25\%}, while the
model still suggested insecure code
\autocite[abstract]{Majdinasab2024Replication}.
In real GitHub projects, 29.5\,\% of the Python and 24.2\,\% of the
JavaScript snippets attributed to AI tools had a security weakness
\autocite[abstract]{Fu2025Weaknesses}, and in the commits of the 1,000 most
popular repositories some weakness families were over-represented in code
labelled as AI-generated \autocite{Wang2025Wild}.
A vendor's preprint reports vulnerabilities in 55.8\,\% of 3,500 generated
artefacts \autocite{Blain2026Broken}, and a systematisation finds that
models recognise security principles that they then fail to apply
\autocite{Patir2026SoK}.
Other field data are milder: of 7,703 files on GitHub explicitly attributed
to AI tools, \enquote{87.9\% of AI-generated code does not contain
identifiable CWE-mapped vulnerabilities} \autocite[abstract]{Schreiber2025GitHub}.
Comparisons with human-written code qualify these rates.
Given the contexts in which developers had introduced C/C++
vulnerabilities, Copilot reproduced the vulnerable code about 33\,\% and the
fix about 25\,\% of the time, and the authors conclude that it \enquote{is
not as bad as human developers at introducing vulnerabilities in code}
\autocite[abstract]{Asare2023Copilot}.
ChatGPT's answers to 108 security-related Java questions had 20\,\% fewer
vulnerabilities than the Stack Overflow answers to the same questions
\autocite[abstract]{Hamer2024Copy}; code that developers obtained from
ChatGPT in real conversations was \enquote{about as likely to contain
vulnerabilities as developer-written code}
\autocite[abstract]{Belozerov2026Repair}; on 72 Python tasks both human and
GPT-4 code had flaws, the human code more varied ones and the GPT-4 code
more severe outliers \autocite{Licorish2025Jury}; and human-written Python
on GitHub had a higher prevalence of security weaknesses than AI-attributed
Python \autocite{CaballeroIgnacio2026Trends} (abstract only), while
language models used security-sensitive functions more often than humans
but less dangerously \autocite{Lee2026HumanAI} (abstract only).

\paragraph{Newer models and agents.}
Across 32 models from seven families, \enquote{newer models do become safer
in absolute terms, although no family closes the gap} between code that
passes functional tests and code that also passes security tests
\autocite[abstract]{SantosdeMoura2026Newer}.
An analysis of four versions each of three model families found marginal
improvement for proprietary models, stagnation for open-weight ones and
occasional regressions \autocite{Cannavale2026Smarter} (abstract only), and
in a 2023 benchmark \enquote{models that are more capable at coding tend to
be more prone to insecure code suggestions}
\autocite[sect.~1]{Bhatt2023CyberSecEval}.
For coding agents on 186 real-world feature tasks, \enquote{although 57\% of
the solutions from SWE-Agent with Claude 4 Sonnet are functionally correct,
only 11.8\% are secure} \autocite[abstract]{Zhao2026VibeCoding}, and agents
copy vulnerabilities that already exist in the code base
\autocite{Sun2026Propagate}.
In real repositories the picture is mixed: compared with 6,620 human pull
requests, 7,200 agent pull requests introduced fewer or less severe issues
on small bug fixes and differed by task otherwise
\autocite{DectotLeMonnier2026Agentic} (abstract only); most agents' patches
for 500 real issues avoided new vulnerabilities \autocite{ChenJiang2025Agents};
and about 4\,\% of 33,000 agent-authored pull requests were security-related,
mostly hardening rather than fixes \autocite{Siddiq2026Agentic}.

\paragraph{Package hallucination.}
Across 16 models and 576,000 generated samples, \enquote{the average
percentage of hallucinated packages is at least 5.2\% for commercial models
and 21.7\% for open-source models}, with 205,474 distinct invented names
\autocite[abstract]{Spracklen2025Package}; the rate depends on the model,
the language, the model's size and the task \autocite{Krishna2025Phantoms}.
Earlier evaluations counted standard-library modules as hallucinations,
which inflated the rate for Python by up to 9.4 percentage points
\autocite{Djire2026Defenses}.
A replication on five 2026 frontier models found rates of 3.6 to 5.0\,\% of
package mentions; about 1\,\% of the invented names were registered within
six months, and one was registered with malware
\autocite{Churilov2026Range}.
A census of names that had circulated publicly found 22 of 149 registered,
of which \enquote{zero are malicious} \autocite{Okwor2026Registration}.
Both are single-author preprints.

\section{Discussion and limitations}
\label{insecure-code:discussion}

The two parts of the question have different answers.
That generated code often contains vulnerabilities is reported by every
measurement found, with different models, benchmarks and field data, but
the rate depends heavily on how the tasks were chosen: the scenarios of
\textcite{Pearce2022Asleep} and the tasks of \textcite{Zhao2026VibeCoding}
were selected because a vulnerability was possible, so their rates are not
base rates for everyday code, and the field data range from a quarter of
snippets \autocite{Fu2025Weaknesses} to an eighth of files
\autocite{Schreiber2025GitHub}.
Whether generated code is worse than code humans write is not established:
the comparisons with human-written code found it similar or better
\autocite{Asare2023Copilot,Hamer2024Copy,Belozerov2026Repair}.
That people with an assistant write less secure code than people without one
rests on one experiment \autocite{Perry2023Users}; five other comparisons
found a small difference, no significant difference, or a benefit on one
task \autocite{Sandoval2023LostC,Asare2024UserCentered,
ZareShahneh2025Copilot,Jost2026Gemini,Yardim2026Password}.
The over-confidence half is weaker still: one study found it, one found
calibrated confidence \autocite{Khalid2026Trust}, and the rest measured trust
in manipulated settings or by self-report.
Most participants in the comparisons were students or freelancers, and the
models were those of 2021 to 2024; no comparison of people with and without
a current coding agent was found.

Limitations of the data and tools: DBLP could not be searched; IEEE Xplore,
Semantic Scholar and arXiv failed for 9, 4 and 9 of the 18 queries, so their
columns undercount; Scopus held no abstracts for the first searches;
OpenAlex searches only the part of an abstract it holds.
Three studies of generated applications that the check by meaning found
describe security in words outside the queries' vulnerability block and are
outside the session.
Seven sources were read in abstract only.
Several sources are preprints \autocite{Jost2026Gemini,Khalid2026Trust,
SantosdeMoura2026Newer,Sun2026Propagate,Churilov2026Range,
Okwor2026Registration,Bhatt2023CyberSecEval}, and one is a vendor's
evaluation of its own tool \autocite{Blain2026Broken}.
None of the databases searched is an education database.
Only the first tier (2,644 of 21,436 distinct records) was screened in full;
the other records were only counted, apart from two seeded random samples of
60.
Of the general sample, 2 bore on the claim (3.3\,\%; 95\,\% modified Wilson
interval 0.6 to 11.4\,\%), and of the sample of counter-query records 1 did
(1.7\,\%; 0.1 to 8.9\,\%).
Since these shares are not near zero, hundreds of the unscreened records may
bear on the question, and the search is therefore not exhaustive: a study
that would change the conclusion could lie among them.
The screen was done by a language model and an AI agent, and the author has
not reviewed any decision.

\section{Conclusion}
\label{insecure-code:conclusion}

Partly: code generated by language models and coding agents often contains
vulnerabilities, and newer models are safer without having closed the gap,
but whether people who program with an assistant write less secure code
than people without one, and are more confident in it, is not settled: one
experiment with mostly students and a 2022 model found both, and later and
larger comparisons found little or no difference.
The notes may therefore say that AI-generated code must be checked for
security, and that one well-known study found its users both less secure
and more confident; they should not present that study's finding as the
general result, nor a single percentage as the current rate of vulnerable
generated code, nor claim that generated code is less secure than code
written by people.
Package hallucination is real and persists in 2026 models, but its
exploitation has so far been observed only rarely.

\section{Future work}
\label{insecure-code:future}

This search left the following undone.
IEEE Xplore, Semantic Scholar and arXiv must be searched again for the
queries that failed, and DBLP when it answers automated clients.
Seven sources were read only in their abstracts and should be read in full:
the password-storage study \autocite{Yardim2026Password}, the longitudinal
analysis \autocite{Cannavale2026Smarter}, the education study
\autocite{Yuan2025Mindset}, the manipulated-assistant study
\autocite{Serafini2025Manipulated}, the comparison of security-sensitive
functions \autocite{Lee2026HumanAI}, the study of agents' pull requests
\autocite{DectotLeMonnier2026Agentic} and the comparison of AI-attributed
and human code on GitHub \autocite{CaballeroIgnacio2026Trends}.
The records outside the first tier were screened only in two samples, and
the author has not reviewed the machine's decisions; a seeded sample of
them should be reviewed and its disagreement rate reported here.

The field lacks a controlled comparison of people with and without a current
coding agent, with participants like the readers of these notes (beginners
outside computer science), large enough to detect a moderate difference, and
with security and confidence both measured.
If such a study found more vulnerabilities and higher confidence with the
agent, the notes could state the claim without reservation; if it found
none, the warning would rest on the tools alone.

\section{Hits that bear on the claim}
\label{insecure-code:hits}

% Classified hit list, written by:
%   scholar sessions export insecure-code -f table --bearing-only --lang en
%     --label insecure-code --track "Insecure code with AI assistants"
%     --theme cited-user-study=... --theme cited-confidence=...
%     --theme cited-prevalence=... --theme cited-model-change=...
%     --theme cited-agents=... --theme cited-package-halluc=...
%     -o litteratursokning/insecure-code-full
% Raw export beside it: litteratursokning/insecure-code-export.csv.
% Rows the agent kept with a cited-* tag show as candidates until the
% author confirms them.
\Cref{tab:insecure-code} lists the records that bear on the claim.
Of the 21,436 distinct records, 2,644 (the first tier) were screened in
full; the others were only counted, apart from two seeded random samples of
60 each, of which 2 and 1 records bore on the claim
(\cref{insecure-code:method}).
The records the table counts as unclassified are those that were counted
but not screened.
% --- scholar LaTeX fragment --------------------------------------------
% Generated by: scholar sessions export -f table
% Session: insecure-code
% \input this file from the body of a document whose
% preamble loads:
%   \usepackage{longtable}
%   \usepackage{booktabs}
%   \usepackage{array}
% ----------------------------------------------------------------------

\begingroup\footnotesize
\begin{longtable}{@{}%
  >{\raggedright\arraybackslash}p{0.44\textwidth}c%
  >{\raggedright\arraybackslash}p{0.13\textwidth}%
  >{\raggedright\arraybackslash}p{0.26\textwidth}@{}}
\caption{Records bearing on the claim, Insecure code with AI assistants (0 cited, 238 supporting and 63 qualifying or contradicting it, out of 21,435 unique hits; the 1,035 adjacent hits, the 1,281 false hits and 18,818 unclassified are not listed). The last column gives the reason for each inclusion; for machine-classified rows the model's confidence is given, and for rows decided by an agent the agent's. Databases searched: OpenAlex (OA; not open access), IEEE Xplore (IEEE; failed for 9 of 18 queries), Web of Science (WoS), Scopus, Semantic Scholar (S2; failed for 4 of 18 queries), arXiv (failed for 8 of 18 queries). Searched 2026-11-27.}\label{tab:insecure-code}\\
\toprule
Title & Year & Database & Reason \\
\midrule
\endfirsthead
\multicolumn{4}{@{}l}{\emph{\tablename~\thetable{} (cont.)}}\\
\toprule
Title & Year & Database & Reason \\
\midrule
\endhead
\midrule \multicolumn{4}{r@{}}{\emph{continued on next page}}\\
\endfoot
\bottomrule
\endlastfoot
\textquotedbl{}Tab, Tab, Bug\textquotedbl{}: Security Pitfalls of Next Edit Suggestions in AI-Integrated IDEs & 2026 & OA, arXiv, S2 & supports the claim (0.70) \\
\textquotedbl{}You still have to study\textquotedbl{} -- On the Security of LLM generated code & 2024 & OA, S2, arXiv & supports the claim (0.85) \\
A Comparative Security Analysis of AI-Generated and Human-Written Code Using Static Vulnerability Detection & 2026 & IEEE, OA, S2 & supports the claim (0.80) \\
A Comparative Security Evaluation of AI-Generated Code Across Large Language Models and Programming Languages & 2026 & OA & supports the claim (0.92) \\
A Pilot Study on Secure Code Generation with ChatGPT for Web Applications & 2024 & WoS, OA, Scopus, S2 & supports the claim (0.55) \\
A systematic literature review on the impact of AI models on the security of code generation & 2024 & WoS, OA, Scopus, S2 & supports the claim (0.80) \\
A.S.E 2.0: Understanding How Repository Context Impacts the Security of LLM-Generated Code & 2026 & OA, arXiv, S2 & supports the claim (0.75) \\
A.S.E: A Repository-Level Benchmark for Evaluating Security in AI-Generated Code & 2025 & S2, arXiv & supports the claim (0.85) \\
Adversarial Testing of Automated Program Repair Agents for Security Vulnerabilities & 2026 & OA, S2, arXiv & supports the claim (0.70) \\
AI ASSISTANTS IN SOFTWARE DEVELOPMENT: ANALYSIS OF SECURITY RISKS IN GENERATED CODE & 2026 & OA, S2 & supports the claim (0.60) (no-evidence) \\
AI Code in the Wild: Measuring Security Risks and Ecosystem Shifts of AI-Generated Code in Modern Software & 2025 & OA, S2, arXiv & supports the claim (0.75) \\
AI Code Synthesis, Autonomous Debugging, and the Evolution of the Software Engineering Lifecycle: Repository Graph Reasoning, SWE-bench Benchmarks, and Security Bounds & 2026 & OA & supports the claim (0.60) \\
AI Writes, We Analyze: The ChatGPT Python Code Saga & 2024 & IEEE, WoS, OA, Scopus & supports the claim (0.70) \\
AI-Based Detection of Hallucinated and Vulnerable Code Generated by LLMs & 2026 & OA & supports the claim (0.60) (new-technique) \\
AI-Induced Supply-Chain Compromise: A Systematic Review of Package Hallucinations and Slopsquatting Attacks & 2025 & OA & supports the claim (0.80) \\
AI-Powered Code Generation in Agile Workflows: Risks, Review Strategies, and Developer Perceptions & 2026 & IEEE, OA, Scopus, S2 & supports the claim (0.55) \\
AIALIB: A Threat Library for AI-Generated Code & 2026 & OA, Scopus, S2 & supports the claim (0.60) \\
An Empirical Evaluation of Default Security Vulnerabilities in AI-Generated Code Across C++ and Java & 2026 & OA & supports the claim (0.70) \\
An Empirical Study of Security Calibration in Large Language Models for Code & 2026 & arXiv, IEEE, OA, S2 & supports the claim (0.75) \\
An empirical study of security vulnerabilities introduced by AI-generated code across programming languages & 2023 & OA, S2 & supports the claim (0.75) \\
An Exploratory Study on Fine-Tuning Large Language Models for Secure Code Generation & 2024 & arXiv, WoS, OA, Scopus, S2 & supports the claim (0.50) (new-technique) \\
An Investigation into Misuse of Java Security APIs by Large Language Models & 2024 & arXiv, WoS, OA, Scopus, S2 & supports the claim (0.90) \\
Análisis de Vulnerabilidades en Vibe Coding para Aplicaciones Web: Un Estudio de Caso & 2026 & OA & supports the claim (0.80) \\
APILOT: Improving the Security and Usability of LLM Code Suggestions via Outdated API Mitigation & 2025 & WoS, IEEE, OA, Scopus, S2 & supports the claim (0.50) (new-technique) \\
Are Large Language Models Safe? A Vulnerability Analysis of Generated Source Code & 2026 & OA & supports the claim (0.80) \\
Are We There Yet? On Security Vulnerabilities Produced by Open Source Generative AI Models and Its Implications for Security Education & 2025 & WoS, Scopus, S2 & supports the claim (0.80) \\
Artifact for: Package hallucination and the extractor's character class & 2026 & OA & supports the claim (0.55) \\
Artificially Insecure: Examining GitHub Copilot's AI-based Vulnerability Prevention System & 2025 & WoS, IEEE, OA, Scopus, S2 & supports the claim (0.90) \\
Asleep at the Keyboard? Assessing the Security of GitHub Copilot's Code Contributions & 2025 & WoS, OA, S2, Scopus, IEEE & supports the claim (0.95) \\
Assessing the Security of ``Vibe Coding'': An Analysis of Vulnerabilities in Web Applications Generated by LLMs & 2026 & OA, S2 & supports the claim (0.85) \\
Assessing the Security of Vibe Coding: Baseline Vs. Security-Oriented Prompts in LLM Code Generation & 2026 & OA, Scopus & supports the claim (0.80) \\
AutoBaxBuilder: Bootstrapping Code Security Benchmarking & 2025 & arXiv, OA, Scopus, S2 & supports the claim (0.80) \\
AutoSUIT Bench - Automated Security UnIt Test Benchmark for LLM Coding & 2026 & OA, Scopus, S2 & supports the claim (0.75) \\
Aware but Unprepared: Measuring the Security Awareness-Behavior Gap in Student Use of LLM-Generated Code with Bifröst & 2025 & OA, arXiv, S2 & supports the claim (0.80) \\
Balancing Security and Correctness in Code Generation: An Empirical Study on Commercial Large Language Models & 2025 & WoS, IEEE, OA, Scopus, S2 & supports the claim (0.70) \\
Bayesian-Calibrated Detection of Hallucinated Package Imports in AI-Assisted Code & 2026 & OA, arXiv, S2 & supports the claim (0.60) (new-technique) \\
Broken by Default: A Formal Verification Study of Security Vulnerabilities in AI-Generated Code & 2026 & OA, arXiv & supports the claim (0.60) \\
Can ChatGPT replace StackOverflow? A Study on Robustness and Reliability of Large Language Model Code Generation & 2023 & arXiv & supports the claim (0.60) \\
Can We Trust Large Language Models Generated Code? A Framework for In-Context Learning, Security Patterns, and Code Evaluations Across Diverse LLMs & 2024 & OA, S2, arXiv & supports the claim (0.70) \\
Can You Really Trust Code Copilot? Evaluating Large Language Models from a Code Security Perspective & 2025 & OA & supports the claim (0.90) \\
Can You Really Trust Code Copilots? Evaluating Large Language Models from a Code Security Perspective & 2025 & OA, Scopus, S2, arXiv & supports the claim (0.80) \\
Can You Really Trust Code Copilots? Evaluating Large Language Models from a Code Security Perspective & 2025 & WoS & supports the claim (0.85) \\
Code Generation of Smart Contracts with LLMs: A Case Study on Hyperledger Fabric & 2025 & WoS, IEEE, OA, Scopus, S2 & supports the claim (0.70) \\
Code Generation Security in LLMs: A Hybrid Detection and Post-Processing Framework for Vulnerability Mitigation & 2025 & OA, S2 & supports the claim (0.60) (new-technique) \\
Code generation with large language models: a survey from neural program synthesis to autonomous software development & 2026 & WoS, OA, Scopus, S2 & supports the claim (0.50) (no-evidence) \\
CodeLMSec Benchmark: Systematically Evaluating and Finding Security Vulnerabilities in Black-Box Code Language Models & 2023 & arXiv, WoS, IEEE, OA, Scopus, S2 & supports the claim (0.70) \\
Coding Agents Propagate Vulnerabilities in Code Bases & 2026 & OA, arXiv, S2 & supports the claim (0.80) \\
Comparative Analysis of Security Vulnerabilities in Ai-Generated Versus Human-Written Python Code: an Automated Static Analysis Approach & 2026 & IEEE, OA, Scopus, S2 & supports the claim (0.85) \\
Context-Based Adversarial Attacks on AI Code Generators: Vulnerability Analysis and Implications & 2026 & OA, arXiv & supports the claim (0.60) \\
Copilot Impact Studies: Measuring Productivity, Trust, and Skill Evolution in Enterprise Developer Teams & 2025 & OA, S2 & supports the claim (0.50) (no-evidence) \\
Cross-Cutting Security Analysis of LLM-Generated Code via Metamorphic Testing and Association Rule Mining & 2026 & arXiv, IEEE, OA, Scopus, S2 & supports the claim (0.80) \\
CWEval: Outcome-driven Evaluation on Functionality and Security of LLM Code Generation & 2025 & arXiv, WoS, IEEE, OA, Scopus, S2 & supports the claim (0.80) \\
Cybersecurity Risks and Software Vulnerabilities in AI-Generated Code & 2026 & OA & supports the claim (0.75) \\
Dataset for the Paper: Where Did the Repair First Go Wrong? Localizing the Origins of Silent Failures in Agentic Vulnerability Repair & 2026 & OA & supports the claim (0.50) \\
De-identified Reproducibility Archive for Beyond Passive Slopsquatting: From Historical Package Hallucinations to Active Dependency Steering & 2026 & OA & supports the claim (0.60) \\
Debt Behind the AI Boom: A Large-Scale Empirical Study of AI-Generated Code in the Wild & 2026 & OA, S2, arXiv & supports the claim (0.70) \\
Do AI Coding Assistants Check Before They Install? A Pre-Registered Demand-Side Audit of Trust Signals in the Research Software Supply Chain & 2026 & arXiv & supports the claim (0.60) \\
Do Coding Agents Catch Their Own Bugs? (They Don't.) & 2026 & OA, S2, arXiv & supports the claim (0.55) \\
Do Language Models Prefer Vulnerable Code? A Probabilistic Study of Insecure Code Preference & 2026 & WoS, IEEE, OA, Scopus, S2 & supports the claim (0.65) \\
Do Users Write More Insecure Code with AI Assistants? & 2023 & WoS, OA, Scopus, arXiv, S2 & supports the claim (0.90) \\
Empirical Measurement of LLM Package Hallucination Risk Across Python, Rust, and Go Ecosystems & 2026 & OA & supports the claim (0.75) \\
Enhancing Large Language Models for Secure Code Generation: A Dataset-driven Study on Vulnerability Mitigation & 2023 & OA, S2, arXiv & supports the claim (0.75) (new-technique) \\
Enhancing Security in Industrial Application Development: Case Study on Self-Generating Artificial Intelligence Tools & 2024 & WoS, OA, Scopus, S2 & supports the claim (0.60) \\
Eradicating the Unseen: Detecting, Exploiting, and Remediating a Path Traversal Vulnerability across GitHub & 2025 & WoS, Scopus, S2 & supports the claim (0.50) \\
Evaluating and Enhancing Memory Safety in LLM-Generated Code on CHERI & 2026 & IEEE, OA, S2 & supports the claim (0.85) (new-technique) \\
Evaluating Large Language Models for Code Generation: Reliability, Security, and Hallucination Risks --- A Systematic Review and Comparative Analysis of Published Benchmark Evidence & 2026 & OA & supports the claim (0.80) (no-evidence) \\
Evaluating Large Language Models in Cybersecurity: A Systematic Taxonomy and Empirical Analysis & 2026 & WoS, OA, Scopus, S2 & supports the claim (0.45) (no-evidence) \\
Evaluating Static Security Gates and Learned Detectors on LLM-Generated Code in CI/CD Pipelines & 2026 & WoS, IEEE, Scopus, OA, S2 & supports the claim (0.70) \\
Evaluating the Effectiveness of LLM Agents Generating API Implementations from Gherkin Scenarios: A Pilot Study & 2026 & OA, Scopus, S2 & supports the claim (0.60) \\
Evaluating the Efficacy of Large Language Models for Automated Vulnerable Code Generation & 2026 & WoS, OA, Scopus & supports the claim (0.60) \\
Evaluating the Impact of Generative AI on Intelligent Programming Assistance and Code Quality. & 2025 & OA, S2 & supports the claim (0.45) \\
Evaluating the Reliability and Security of AI-Generated Source Code: A Comparative Study of Large Language Models & 2026 & OA, S2 & supports the claim (0.40) \\
Evaluating the Vulnerability Landscape of LLM-Generated Smart Contracts & 2026 & OA, S2, arXiv & supports the claim (0.85) \\
Explicit Vulnerability Generation with LLMs: An Investigation beyond Adversarial Attacks & 2025 & WoS, IEEE, Scopus, OA, S2, arXiv & supports the claim (0.60) \\
Exploring the Efficacy of Large Language Models in Full-stack Development : A Study on the Quality of ChatGPT Generated Code in Full-stack Development & 2026 & OA & supports the claim (0.70) \\
Exploring the Impact of Intervention Methods on Developers' Security Behavior in a Manipulated ChatGPT Study & 2025 & WoS, OA, Scopus, S2 & supports the claim (0.70) \\
Exploring the Security Landscape of AI Generated Code: A Comprehensive Study & 2026 & IEEE, OA, Scopus, S2 & supports the claim (0.50) (no-evidence) \\
Exploring the Security Threats of Knowledge Base Poisoning in Retrieval-Augmented Code Generation & 2025 & arXiv, WoS, IEEE, OA, Scopus, S2 & supports the claim (0.70) \\
False Security Confidence in Benign LLM Code Generation & 2026 & OA, S2, arXiv & supports the claim (0.50) (no-evidence) \\
febinrenu/slopsquat-study: v1.0.0: Hallucinated package names attract malware & 2026 & OA & supports the claim (0.80) \\
Fight Fire with Fire: How Much Can We Trust ChatGPT on Source Code-Related Tasks? & 2024 & arXiv, WoS, IEEE, OA, Scopus, S2 & supports the claim (0.70) \\
Fine Tuning Large Language Model for Secure Code Generation & 2024 & WoS, IEEE, OA, Scopus, S2 & supports the claim (0.45) (new-technique) \\
From Code Synthesis to Security Analysis: Evaluation of Open-Source LLMs on Ethereum Smart Contracts & 2026 & IEEE, OA, Scopus, S2 & supports the claim (0.60) \\
From Model to Breach: Towards Actionable LLM-Generated Vulnerabilities Reporting & 2025 & OA, S2, arXiv & supports the claim (0.80) \\
From Preventive to Reactive: How AI Coding Assistants Transform Developers' Security Awareness & 2026 & WoS, OA, arXiv, S2 & supports the claim (0.70) \\
From Vulnerabilities to Remediation: A Systematic Literature Review of LLMs in Code Security & 2024 & S2, OA, arXiv & supports the claim (0.70) \\
Generative AI and Large Language Models for Software Engineering: Opportunities, Challenges, and Future Directions & 2026 & OA & supports the claim (0.80) \\
Guiding AI to Fix Its Own Flaws: An Empirical Study on LLM-Driven Secure Code Generation & 2026 & OA, S2, arXiv & supports the claim (0.70) \\
Hal-Detector: Is Explainable AI Helping Detect Slopsquatting? & 2026 & OA, S2 & supports the claim (0.50) (new-technique) \\
HalluGuard: A Multi-Signal Detection Framework for Hallucinations in LLM-Generated Code & 2026 & IEEE, OA, Scopus, S2 & supports the claim (0.65) \\
Harmless Yet Harmful: Neutral Prompting Attacks for Stealthy Hallucination Steering in Agent Skills & 2026 & OA, arXiv, S2 & supports the claim (0.60) \\
Helping LLMs improve code generation using feedback from testing and static analysis & 2026 & OA, Scopus, S2, arXiv & supports the claim (0.70) \\
HFuzzer: Testing Large Language Models for Package Hallucinations via Phrase-based Fuzzing & 2025 & WoS, OA, Scopus, arXiv, S2 & supports the claim (0.75) (new-technique) \\
How Developers Respond to Insecure AI Code Suggestions under Poisoning Attacks: An Empirical Study of Interface-Level Countermeasures & 2026 & IEEE, OA, Scopus & supports the claim (0.70) \\
How Safe Are AI-Generated Patches? A Large-scale Study on Security Risks in LLM and Agentic Automated Program Repair on SWE-bench & 2025 & arXiv, OA, S2 & supports the claim (0.85) \\
How Secure are Full-Stack Web Applications Generated by Large Language Models? A Multi-language Security Study & 2027 & OA, Scopus & supports the claim (0.90) \\
How secure is AI-generated Code: A Large-Scale Comparison of Large Language Models & 2024 & arXiv, WoS, OA, Scopus, S2 & supports the claim (0.92) \\
How Secure is Code Generated by ChatGPT? & 2023 & arXiv, IEEE, OA, Scopus, S2 & supports the claim (0.90) \\
How Well Do Large Language Models Serve as End-to-End Secure Code Agents for Python? & 2025 & WoS, OA, Scopus, S2, arXiv & supports the claim (0.90) \\
How Well Do Large Language Models Serve as End-to-End Secure Code Producers? & 2024 & S2, arXiv & supports the claim (0.90) \\
Human-Written vs. AI-Generated Code: A Large-Scale Study of Defects, Vulnerabilities, and Complexity & 2025 & arXiv, WoS, IEEE, OA, Scopus, S2 & supports the claim (0.90) \\
Impact of Large Language Models (LLMs) and Generative AI on Backend Coding: A Systematic Literature Review & 2026 & OA, Scopus & supports the claim (0.70) \\
Implementation of Static Code Analysis to Detect Vulnerabilities in Applications Developed with the Assistance of Large-Language Models (LLM) & 2025 & OA, S2 & supports the claim (0.70) \\
Importing Phantoms: Measuring LLM Package Hallucination Vulnerabilities & 2025 & OA, S2, arXiv & supports the claim (0.80) \\
ImportSnare: Directed \textquotedbl{}Code Manual\textquotedbl{} Hijacking in Retrieval-Augmented Code Generation & 2025 & arXiv, WoS, OA, Scopus, S2 & supports the claim (0.55) \\
In-Context Vulnerability Propagation in LLMs [Work In Progress Paper] & 2025 & WoS, Scopus, OA, S2 & supports the claim (0.60) \\
Inclusive or Exploitable? Cultural Exclusion and ReDoS Vulnerabilities in LLM-Generated Input Validation Code --- Replication package & 2026 & OA & supports the claim (0.55) \\
Insecure by design? A human-centric security perspective on AI-assisted software development & 2026 & WoS, OA, Scopus, S2 & supports the claim (0.40) (no-evidence) \\
Insecure Output Handling in Large Language Models / Niebezpieczne przetwarzanie danych wyjściowych w dużych modelach językowych & 2025 & OA & supports the claim (0.70) \\
Insights into Security-Related AI-Generated Pull Requests & 2026 & OA, S2, arXiv & supports the claim (0.80) \\
INTEGRATION OF ARTIFICIAL INTELLIGENCE INTO THE SOFTWARE DEVELOPMENT LIFE CYCLE (SDLC) & 2026 & OA, S2 & supports the claim (0.80) \\
Is AI-Generated Web Code Vulnerability-Free? & 2026 & IEEE, OA, Scopus, S2 & supports the claim (0.85) \\
Is Vibe Coding Safe? Benchmarking Vulnerability of Agent-Generated Code in Real-World Tasks & 2026 & OA, Scopus, S2, arXiv & supports the claim (0.90) \\
Just Pin It: Can Coding Agents Be Trusted to Choose Dependency Versions? & 2026 & OA & supports the claim (0.70) \\
Large Language Models and Simple, Stupid Bugs & 2023 & WoS, IEEE, OA, Scopus, S2 & supports the claim (0.50) \\
Large Language Models for Code: Security Hardening and Adversarial Testing & 2023 & WoS, OA, Scopus, S2, arXiv & supports the claim (0.60) \\
Last Secure Code Benchmark: Is the Code You Vibe Secure? & 2026 & OA, arXiv, S2 & supports the claim (0.95) \\
Library Hallucinations in LLM-Generated Code: A Risk Analysis Grounded in Developer Queries & 2025 & arXiv, S2 & supports the claim (0.85) \\
LLM Ghostbusters: Surgical Package Hallucination Suppression via Adaptive Unlearning & 2026 & OA, S2, arXiv & supports the claim (0.60) (new-technique) \\
LLM-CSEC: Empirical Evaluation of Security in C/C++ Code Generated by Large Language Models & 2025 & arXiv, OA, WoS, Scopus, S2 & supports the claim (0.88) \\
LLMs in Web Development: Evaluating LLM-Generated PHP Code Unveiling Vulnerabilities and Limitations & 2024 & WoS, OA, Scopus, arXiv, S2 & supports the claim (0.90) \\
LLMSecEval: A Dataset of Natural Language Prompts for Security Evaluations & 2023 & arXiv, WoS, IEEE, OA, Scopus, S2 & supports the claim (0.65) \\
Minimal Prompt Perturbations Lead to Code Vulnerabilities: Prompt Fragility and Hidden-State Signals in Coding LLMs & 2026 & OA, S2, arXiv & supports the claim (0.80) \\
Mitigating Package Hallucinations in Large Language Models via Model Editing & 2026 & OA, arXiv, S2 & supports the claim (0.60) (new-technique) \\
MOSAIC-Bench: Measuring Compositional Vulnerability Induction in Coding Agents & 2026 & OA, S2, arXiv & supports the claim (0.90) \\
Names Can Hurt: Spotting Slopsquatting Risks Caused by Package Name Hallucinations in Local Coding LLMs & 2026 & OA, arXiv, S2 & supports the claim (0.50) (new-technique) \\
NatSecBench: Measuring Vulnerability Injection Rates in Deployed LLM Code Assistants Under Naturalistic Developer Prompts & 2026 & OA & supports the claim (0.85) \\
Navigating (in)Security of AI-Generated Code & 2024 & WoS, IEEE, OA, Scopus, S2 & supports the claim (0.90) \\
No Need to Lift a Finger Anymore? Assessing the Quality of Code Generation by ChatGPT & 2023 & arXiv, WoS, IEEE, OA, Scopus, S2 & supports the claim (0.85) \\
On the Gap Between Declarative Vulnerability Knowledge and Secure Code Generation in Large Language Models & 2026 & OA & supports the claim (0.80) \\
Package Hallucination Attacks on Coding Agents through Prompt Injection in Rule Files & 2026 & OA, S2, arXiv & supports the claim (0.50) (attack-on-model) \\
Package hallucination of quantised local code models and a registry-grounded decoding guard: code, prompts, redacted outputs and results & 2026 & OA & supports the claim (0.65) \\
Package Hallucinations as Phantoms in Open Source Software Supply Chains: An Empirical Security Analysis & 2026 & OA, S2 & supports the claim (0.90) \\
PackMonitor: Enabling Zero Package Hallucinations Through Decoding-Time Monitoring & 2026 & arXiv, S2 & supports the claim (0.70) (new-technique) \\
PackMonitor: Towards Zero Package Hallucinations via Decoding-Time Monitoring & 2026 & OA, S2 & supports the claim (0.60) (new-technique) \\
Perceived versus Actual Security in Web Applications : A Comparative Analysis of Human-Developed and Large Language Model (LLM)-Generated Applications & 2026 & OA & supports the claim (0.70) \\
Persistent Human Feedback, LLMs, and Static Analyzers for Secure Code Generation and Vulnerability Detection & 2026 & arXiv, WoS, IEEE, OA, Scopus, S2 & supports the claim (0.55) (new-technique) \\
Phantom Dependencies: How and Why Code LLMs Hallucinate Software Packages & 2026 & OA & supports the claim (0.50) \\
Poisoned ChatGPT Finds Work for Idle Hands: Exploring Developers' Coding Practices with Insecure Suggestions from Poisoned AI Models & 2023 & arXiv, WoS, IEEE, OA, Scopus, S2 & supports the claim (0.80) \\
Poster: Insecure Coding Habits Die Hard. Can PEFT Really Turn LLMs into Secure Coders? & 2025 & WoS, OA, Scopus, S2 & supports the claim (0.45) (new-technique) \\
PRECEPT: Security Self-Improvement of Coding Agents via Evolving Principles & 2026 & OA, S2, arXiv & supports the claim (0.70) (new-technique) \\
Premature Trust: How Student Overreliance on GenAI is Seeding Tomorrow's Security Gaps & 2025 & WoS, OA, Scopus, S2 & supports the claim (0.40) (no-evidence) \\
Preventing the Introduction of Security Defects when using AI Coding Assistants & 2026 & OA & supports the claim (0.95) \\
PromSec: Prompt Optimization for Secure Generation of Functional Source Code with Large Language Models (LLMs) & 2024 & WoS, OA, Scopus, S2, arXiv & supports the claim (0.45) (new-technique) \\
Protecting the Whisper: A Security Assessment of Amazon CodeWhisperer's Generated Code & 2025 & WoS, OA, Scopus & supports the claim (0.90) \\
Purple Llama CyberSecEval: A Secure Coding Benchmark for Language Models & 2023 & OA, S2 & supports the claim (0.80) \\
PwnPilot: Reflections on Trusting Trust in the Age of Large Language Models and AI Code Assistants & 2023 & IEEE, OA, Scopus, S2 & supports the claim (0.50) (no-evidence) \\
Q\&AEval: Benchmarking Secure Coding Ability of LLMs on Real-World Tasks & 2026 & WoS, OA, Scopus, S2 & supports the claim (0.90) \\
Quality Assessment of ChatGPT Generated Code and their Use by Developers & 2024 & WoS, IEEE, OA, Scopus, S2 & supports the claim (0.80) \\
RealSec-bench: A Benchmark for Evaluating Secure Code Generation in Real-World Repositories & 2026 & S2, arXiv, OA, Scopus & supports the claim (0.85) \\
Redefining the Programmer: Human-AI Collaboration, LLMs, and Security in Modern Software Engineering & 2025 & WoS, OA, Scopus & supports the claim (0.55) \\
Refusing the Impossible: A Taxonomy and Benchmark for Code Hallucination in Large Language Models & 2026 & OA, S2, arXiv & supports the claim (0.80) \\
Research on Security Vulnerability Patterns of Code Generated by Large Language Models & 2026 & IEEE, OA, Scopus, S2 & supports the claim (0.80) \\
SafeGenBench: A Benchmark Framework for Security Vulnerability Detection in LLM-Generated Code & 2025 & OA, S2, arXiv & supports the claim (0.85) \\
SALLM: Security Assessment of Generated Code & 2024 & WoS, IEEE, OA, Scopus, S2, arXiv & supports the claim (0.70) \\
SecProbe: Adaptive Evaluation of Coding Agents on Cybersecurity Vulnerabilities & 2026 & OA, S2, arXiv & supports the claim (0.45) \\
SecRepoBench: Benchmarking Code Agents for Secure Code Completion in Real-World Repositories & 2026 & OA, Scopus, S2, arXiv & supports the claim (0.90) \\
SecRepoBench: Benchmarking LLMs for Secure Code Generation in Real-World Repositories & 2025 & S2, arXiv & supports the claim (0.85) \\
Secure Code Generation with LLMs: Risk Assessment and Mitigation Strategies & 2025 & OA, S2 & supports the claim (0.75) \\
Secure Code Generation With Open Source Generative AI Models & 2026 & WoS, IEEE, OA, Scopus, S2 & supports the claim (0.70) \\
Secure Coding Drift in LLM-Assisted Post-Quantum Cryptography Development: A Gamified Fix & 2026 & OA, S2, arXiv & supports the claim (0.40) (no-evidence; new-technique) \\
Secure or Suspect? Investigating Package Hallucinations of Shell Command in Original and Quantized LLMs & 2025 & OA, S2, arXiv & supports the claim (0.85) \\
SecureForge: Finding and Preventing Vulnerabilities in LLM-Generated Code via Prompt Optimization & 2026 & OA, S2, arXiv & supports the claim (0.75) (new-technique) \\
SecureVibeBench: Benchmarking Secure Vibe Coding of AI Agents via Reconstructing Vulnerability-Introducing Scenarios & 2025 & arXiv, WoS, OA, Scopus, S2 & supports the claim (0.90) \\
SecureVibeBench: Benchmarking Secure Vibe Coding of AI Agents via Reconstructing Vulnerability-Introducing Scenarios & 2026 & OA & supports the claim (0.85) \\
Securing AI-Generated Code & 2025 & OA & supports the claim (0.45) \\
Securing LLM code generation: Leveraging prompt engineering to mitigate vulnerabilities across models and languages & 2026 & WoS, OA, Scopus, S2 & supports the claim (0.85) \\
Securing LLM-Generated Embedded Firmware Through AI Agent-Driven Validation and Patching & 2025 & IEEE, WoS, OA, Scopus, S2, arXiv & supports the claim (0.45) (new-technique) \\
Security Analysis of Automated Code Generation: Structural Vulnerabilities in AI-Generated & 2025 & WoS, OA, Scopus, S2 & supports the claim (0.90) \\
Security Analysis of LLM-Generated Web API Backends & 2026 & OA, WoS, IEEE, Scopus, S2 & supports the claim (0.85) \\
Security Degradation in Iterative AI Code Generation -- A Systematic Analysis of the Paradox & 2025 & arXiv, WoS, IEEE, OA, Scopus, S2 & supports the claim (0.93) \\
Security of LLM-generated Code: A Comparative Analysis & 2026 & OA, S2, arXiv & supports the claim (0.90) \\
Security Risks in AI-Generated Code: A Comparative Analysis of Industry and Academic Benchmarks & 2025 & IEEE, OA, Scopus, S2 & supports the claim (0.60) \\
Security Vulnerabilities in AI-Generated JavaScript: A Comparative Study of Large Language Models & 2025 & WoS, IEEE, OA, Scopus, S2 & supports the claim (0.92) \\
Security Vulnerability Analysis Of AI-Generated Code: A Penetration Testing Approach & 2026 & OA & supports the claim (0.70) \\
Security Vulnerability Patterns in AI-Generated Code: A Cross-Model Comparative Study & 2026 & OA, arXiv & supports the claim (0.90) \\
Security Weaknesses in LLM-Generated Source Code: An Empirical Vulnerability Analysis of Iterative AI-Assisted Development & 2026 & OA & supports the claim (0.85) \\
Security Weaknesses in LLM-Generated Source Code: An Empirical Vulnerability Analysis of Iterative AI-Assisted Development v1.0 & 2026 & OA & supports the claim (0.80) \\
Security Weaknesses of Code Generated by Generative AI & 2026 & OA, Scopus & supports the claim (0.90) \\
Security Weaknesses of Copilot Generated Code in GitHub & 2023 & S2, arXiv & supports the claim (0.90) \\
Security Weaknesses of Copilot-Generated Code in GitHub Projects: An Empirical Study & 2025 & WoS, OA, Scopus, S2 & supports the claim (0.90) \\
Security-Related Code Hallucinations in Large Language Models: A Systematic Literature Review & 2026 & IEEE, OA, Scopus, S2 & supports the claim (0.60) (no-evidence) \\
SecurityEval Dataset: Mining Vulnerability Examples to Evaluate Machine Learning-Based Code Generation Techniques & 2022 & WoS, OA, Scopus, S2 & supports the claim (0.80) \\
Shape Before You Build: Secure Cryptographic Code Generation via Prompt Optimization & 2026 & WoS, OA, Scopus & supports the claim (0.55) (new-technique) \\
Slopsquatting and package-hallucination in LLMS & 2026 & OA, S2 & supports the claim (0.85) \\
Smarter, but Not Safer: An Empirical Analysis of the Functional-Security Gap in Evolving LLMs & 2026 & IEEE, OA, S2 & supports the claim (0.75) \\
SoK: AI Secure Code Generation: Progress, Pitfalls, and Paths Forward & 2026 & OA, S2, arXiv & supports the claim (0.75) \\
SOK: Exploring Hallucinations and Security Risks in AI-Assisted Software Development with Insights for LLM Deployment & 2025 & arXiv, WoS, IEEE, Scopus, S2 & supports the claim (0.50) (no-evidence) \\
SolGen: Secure Smart Contract Code Generation Using Large Language Models Via Masked Prompting & 2025 & WoS, OA, Scopus, S2 & supports the claim (0.45) (new-technique) \\
SOSecure: The Wisdom of the Crowd for Safer AI-Generated Code & 2026 & OA, Scopus, S2 & supports the claim (0.55) (new-technique) \\
Static Analysis as a Feedback Loop: Enhancing LLM-Generated Code Beyond Correctness & 2025 & arXiv, WoS, IEEE, OA, Scopus, S2 & supports the claim (0.60) \\
Static Analysis of Vulnerabilities That AI-Based Code Assistants May Introduce & 2026 & IEEE, Scopus, S2 & supports the claim (0.85) \\
SWE-chat: Coding Agent Interactions From Real Users in the Wild & 2026 & OA, S2, arXiv & supports the claim (0.80) \\
Systematic Evaluation of Security Flaws in Low-Level Code Generated by LLMs & 2026 & OA, Scopus, S2 & supports the claim (0.85) \\
Taxonomy of challenges influencing secure code generation using LLMs: An integrated approach of systematic literature review and F-AHP & 2026 & WoS, OA, Scopus, S2 & supports the claim (0.50) \\
The AI Coding Assistant as Attack Surface: Systematic Exploitation of LLM Package Hallucination at Scale & 2026 & OA & supports the claim (0.50) (no-evidence) \\
The Double-Edged Sword of AI Pair Programmers: A Systematic Literature Review of Security Vulnerabilities in AI-Generated Code and Agentic Development Environments & 2026 & OA, S2 & supports the claim (0.90) \\
The FormAI Dataset: Generative AI in Software Security through the Lens of Formal Verification & 2023 & WoS, OA, Scopus & supports the claim (0.90) \\
The Hidden Risks of LLM-Generated Web Application Code: A Security-Centric Evaluation of Code Generation Capabilities in Large Language Models & 2026 & OA, Scopus, S2, arXiv & supports the claim (0.90) \\
The Impact of Security Mindset on the Use of AI Assistants in Computing Education & 2025 & IEEE, WoS, OA, Scopus, S2 & supports the claim (0.70) \\
The Quiet Contributions: Insights into AI-Generated Silent Pull Requests & 2026 & OA, Scopus, S2, arXiv & supports the claim (0.65) \\
The Role of AI in Software Development: Industry Practices, Security Implications, and Developer Perceptions & 2024 & OA & supports the claim (0.75) \\
The Security of LLM-Generated Code & 2026 & OA & supports the claim (0.90) \\
The Trustworthiness of AI-Generated Code: A Survey of How Students and Working Professionals Trust, Check and Use Code Written by Generative AI & 2026 & OA & supports the claim (0.70) \\
The Vibe Shift in Software Engineering: Evaluating AI-Led Conversational Programming for Performance, Cognition, and Responsible Adoption & 2026 & OA, Scopus, arXiv, S2 & supports the claim (0.80) \\
Toward a Taxonomy of Vulnerability Patterns in AI-Assistant-Generated Code & 2026 & OA & supports the claim (0.50) \\
Toward Autonomous AI-Driven Software Development: A Systematic Review of the Empirical Evidence on Agentic Systems (2022--2026) & 2026 & OA & supports the claim (0.60) \\
Toward Secure Code Generation: Bridging Correctness and Security via Task-Adaptive Vulnerability Modeling and Execution-Based Benchmarking & 2024 & arXiv, OA, S2 & supports the claim (0.60) (new-technique) \\
Towards Secure Code Generation With LLMs: A Study on Common Weakness Enumeration & 2025 & WoS, IEEE, OA, Scopus, S2 & supports the claim (0.50) (new-technique) \\
Towards Tracing the Quality of Code Generated by LLM Back to the Training Code & 2027 & WoS, OA, Scopus & supports the claim (0.50) \\
Trust but Verify? Uncovering the Security Debt of Autonomous Coding Agents & 2026 & OA, arXiv, S2 & supports the claim (0.95) \\
Trust, Verification Practices and Security Knowledge When Using AI-Generated Code Among Computing Students and Earl-Career Developers & 2026 & OA & supports the claim (0.40) \\
Trusting the Machine: How Secure is LLM-Generated RTL Code? & 2025 & IEEE, OA, Scopus, S2 & supports the claim (0.85) \\
Trustworthy Integration of Large Language Models in Software Engineering: Risks, Evaluation, and Mitigation & 2026 & OA & supports the claim (0.70) \\
Understanding and Mitigating Library-Related Issues in LLM-Generated Code & 2026 & OA, arXiv, S2 & supports the claim (0.45) (new-technique) \\
Understanding Regular Expression Denial of Service (ReDoS): Insights from LLM-Generated Regexes and Developer Forums & 2024 & WoS, IEEE, OA, Scopus, S2 & supports the claim (0.75) \\
Understanding the (In)Security of Vibe-Coded Applications & 2026 & OA, S2, arXiv & supports the claim (0.95) \\
Understanding the Impact of AI Code Assistants on Security API Usage: An Empirical Study & 2026 & S2, OA, Scopus, arXiv & supports the claim (0.85) \\
Using AI Assistants in Software Development: A Qualitative Study on Security Practices and Concerns & 2024 & WoS, OA, Scopus, arXiv, S2 & supports the claim (0.50) \\
Vibe coding at risk: A CoV-RAG framework for mitigating slopsquatting attacks in AI-generated code & 2026 & OA, S2 & supports the claim (0.50) (new-technique) \\
Vibe Coding in Practice: Motivations, Challenges, and a Future Outlook -- a Grey Literature Review & 2025 & arXiv & supports the claim (0.55) \\
Vibe Coding Vulnerabilities in Cyber Security Defense & 2026 & OA & supports the claim (0.65) \\
Vibe Coding: A New Paradigm in Software Development Assisted by Artificial Intelligence & 2026 & IEEE, OA, Scopus, S2 & supports the claim (0.70) \\
Vibe Coding: Practice, Performance, Productivity, and Risk -A State-of-the-Art Review & 2026 & OA, arXiv, S2 & supports the claim (0.70) \\
Vibe-Coding and Emerging Cybersecurity Vulnerabilities: A Case Study of The Moltbook Data Exposure & 2026 & OA, S2 & supports the claim (0.60) \\
Vulnerability Assessment of LLM-Generated Smart Contracts in Ethereum & 2025 & WoS, Scopus & supports the claim (0.80) \\
Vulnerability Assessment of LLM-Generated Smart Contracts in Ethereum & 2025 & OA & supports the claim (0.80) \\
We Have a Package for You! A Comprehensive Analysis of Package Hallucinations by Code Generating LLMs & 2025 & WoS, OA, Scopus, arXiv, S2 & supports the claim (0.95) \\
Weaknesses in LLM-Generated Code for Embedded Systems Networking & 2024 & WoS, IEEE, OA, Scopus & supports the claim (0.90) \\
When AI Takes the Wheel: Security Analysis of Framework-Constrained Program Generation & 2026 & WoS, OA, Scopus, S2, arXiv & supports the claim (0.95) \\
When Developer Aid Becomes Security Debt: A Systematic Analysis of Insecure Behaviors in LLM Coding Agents & 2025 & OA, S2, arXiv & supports the claim (0.90) \\
When LLMs Invent Rust Crates: An Empirical Study of Hallucination Patterns and Mitigation & 2026 & OA, arXiv, S2 & supports the claim (0.80) \\
When Packages Lie: Mitigating Package Hallucinations with a Sufficient Context RAG Framework & 2027 & OA, Scopus & supports the claim (0.75) (new-technique) \\
Where Did the Repair First Go Wrong? Localizing the Origins of Silent Failures in Agentic Vulnerability Repair & 2026 & OA, S2, arXiv & supports the claim (0.60) \\
Who Wrote This Code? Investigating Human and AI Programming Patterns & 2025 & IEEE, Scopus, S2 & supports the claim (0.40) (no-evidence) \\
Why AI-Generated Code Is Vulnerable: Cause Taxonomy, Threat Mapping, and CI/CD Security Controls & 2026 & OA & supports the claim (0.85) \\
You Still have to Study On the Security of LLM Generated Code & 2025 & WoS, OA, Scopus & supports the claim (0.85) \\
YUL: Bringing Software Supply Chain Security to Coding Agents & 2026 & OA, S2, arXiv & supports the claim (0.60) (new-technique) \\
(Don't) Trust, but (Don't) Verify: Developers' Attention to Security in AI-Generated Code & 2026 & OA, arXiv, S2 & qualifies or contradicts the claim (0.85) \\
\textquotedbl{}The AI tool can't make it any worse.\textquotedbl{} Investigating Developers' Security Behavior with AI Assistants in a Password Storage Study & 2026 & WoS, OA, Scopus, S2 & qualifies or contradicts the claim (0.70) \\
A Comparative Analysis between AI Generated Code and Human Written Code: A Preliminary Study & 2024 & WoS, IEEE, OA, Scopus, S2 & qualifies or contradicts the claim (0.50) \\
A Comparative Study of Human-Written Code and AI-Generated Code & 2026 & OA & qualifies or contradicts the claim (0.35) \\
A Systematic Evaluation of AI-Generated Security Patches: Analyzing and Exploiting AI-Generated Fixes & 2025 & WoS, OA, Scopus, S2 & qualifies or contradicts the claim (0.80) \\
A Systematic Review of AI-Assisted Code Generation in Regulated Software Environments: Reliability, Governance, and Compliance Implications & 2026 & OA & qualifies or contradicts the claim (0.85) \\
A User-centered Security Evaluation of Copilot & 2024 & WoS, IEEE, OA, Scopus, arXiv, S2 & qualifies or contradicts the claim (0.95) \\
AI-Assisted Code Completion in Modern IDEs: A Comparative Study & 2026 & OA & qualifies or contradicts the claim (0.50) \\
An Empirical Evaluation of LLM-Generated Code Security Across Prompting Methods & 2026 & OA, S2, arXiv & qualifies or contradicts the claim (0.80) \\
An Empirical Study of Generative AI for CRUD Code: A PHP Case with Tutorial Based Novice Baselines & 2025 & IEEE, OA, Scopus, S2 & qualifies or contradicts the claim (0.55) \\
Assessing the Effectiveness of ChatGPT in Secure Code Development: A Systematic Literature Review & 2025 & WoS, OA, Scopus, S2 & qualifies or contradicts the claim (0.60) \\
Assessing the Impact of AI Coding Agents on Security Vulnerabilities in Web Development: An Analysis of Emerging Software Engineering Practices in Zimbabwe & 2026 & OA, S2 & qualifies or contradicts the claim (0.90) \\
Assessing the Promise and Pitfalls of ChatGPT for Automated CS1-driven Code Generation & 2024 & Scopus, S2 & qualifies or contradicts the claim (0.50) \\
Assessing the Security of GitHub Copilot Generated Code -- A Targeted Replication Study & 2023 & arXiv & qualifies or contradicts the claim (0.90) \\
Assessing the Security of GitHub Copilot's Generated Code - A Targeted Replication Study & 2024 & WoS, IEEE, OA, Scopus, S2 & qualifies or contradicts the claim (0.85) \\
Benchmarking Prompt Engineering Techniques for Secure Code Generation with GPT Models & 2025 & arXiv, WoS, IEEE, OA, Scopus, S2 & qualifies or contradicts the claim (0.80) \\
Beyond Bypass: Measuring Vulnerability Regression in LLM-Generated Security Patches & 2026 & OA, Scopus, S2 & qualifies or contradicts the claim (0.80) \\
churik5/slopsquatting-replication-2026: v0.2 --- Preprint release (clean repository) & 2026 & OA & qualifies or contradicts the claim (0.70) \\
Comparing Human and LLM Generated Code: The Jury is Still Out! & 2025 & OA, arXiv & qualifies or contradicts the claim (0.80) \\
Deploy, but verify: Analysing LLM Generated Code Safety & 2025 & IEEE, OA, Scopus, S2 & qualifies or contradicts the claim (0.75) \\
Evaluating Capabilities and Perspectives of Generative AI Tools in Smart Contract Development & 2025 & OA, S2 & qualifies or contradicts the claim (0.50) \\
Evaluating Inference-Time Defenses against Package Hallucination in LLM-Generated Code & 2026 & OA, Scopus, arXiv, S2 & qualifies or contradicts the claim (0.80) \\
Evaluating Security and Inclusivity in LLM-Generated Code: A Controlled Experiment & 2026 & OA & qualifies or contradicts the claim (0.60) (no-abstract) \\
Evaluating Software Development Agents: Patch Patterns, Code Quality, and Issue Complexity in Real-World GitHub Scenarios & 2024 & arXiv, WoS, IEEE, OA, Scopus, S2 & qualifies or contradicts the claim (0.80) \\
Evaluating the Code Quality of AI-Assisted Code Generation Tools: An Empirical Study on GitHub Copilot, Amazon CodeWhisperer, and ChatGPT & 2023 & OA, S2, arXiv & qualifies or contradicts the claim (0.50) \\
From Hallucination to Registration: Do the Package Names That LLMs Invent Actually Get Claimed? & 2026 & OA & qualifies or contradicts the claim (0.75) \\
From Vibe Coding to Secure by Design? : An Empirical Evaluation of Security-Oriented Prompt Engineering for LLM-Generated APIs & 2026 & OA & qualifies or contradicts the claim (0.60) \\
GitHub Copilot and Code Security: An Empirical Evaluation in the Context of Developers' Professional Experience & 2026 & OA, Scopus, S2 & qualifies or contradicts the claim (0.90) \\
Helpful or Harmful? Evaluating LLM-Assisted Vulnerability Patching via a Human Study & 2026 & OA, arXiv, S2 & qualifies or contradicts the claim (0.70) \\
How Are We Doing With Using AI-Based Programming Assistants For Privacy-Related Code Generation? The Developers' Experience & 2025 & arXiv, WoS, OA, Scopus, S2 & qualifies or contradicts the claim (0.50) \\
Human vs AI: Insecure Use of Security-Sensitive Functions in AI-generated Code in Comparison to Human-written Code & 2026 & WoS, OA, Scopus, S2 & qualifies or contradicts the claim (0.70) \\
Is GitHub's Copilot as Bad as Humans at Introducing Vulnerabilities in Code? & 2022 & arXiv, WoS, OA, Scopus, S2 & qualifies or contradicts the claim (0.90) \\
Is Your Prompt Poisoning Code? Defect Induction Rates and Security Mitigation Strategies & 2025 & arXiv, WoS, OA, Scopus, S2 & qualifies or contradicts the claim (0.80) \\
Just another copy and paste? Comparing the security vulnerabilities of ChatGPT generated code and StackOverflow answers & 2024 & WoS, IEEE, OA, Scopus, arXiv, S2 & qualifies or contradicts the claim (0.90) \\
LLM-Assisted Porting of Security-Critical C Libraries to Idiomatic Rust: A Multi-Model Empirical Study & 2026 & WoS, OA, Scopus, S2 & qualifies or contradicts the claim (0.60) \\
Lost at C: A User Study on the Security Implications of Large Language Model Code Assistants & 2023 & WoS, OA, Scopus, arXiv, S2 & qualifies or contradicts the claim (0.95) \\
My Code Is Less Secure with Gen AI: Surveying Developers' Perceptions of the Impact of Code Generation Tools on Security & 2025 & WoS, OA, Scopus, S2 & qualifies or contradicts the claim (0.80) \\
Newer and Bigger, but Safer? A Longitudinal Study of the Functionality-Security Gap in LLM-Generated Code & 2026 & OA, arXiv, S2 & qualifies or contradicts the claim (0.90) \\
On Fixing Insecure AI-Generated Code through Model Fine-Tuning and Prompting Strategies & 2026 & OA, S2, arXiv & qualifies or contradicts the claim (0.60) \\
Open Data Repository for \textquotedbl{}The Limitations of Prompt Engineering in Mitigating LLM Code Vulnerabilities\textquotedbl{} (50 Primary Studies) & 2026 & OA & qualifies or contradicts the claim (0.60) \\
Package Hallucination in LLM-Generated Python Code: Data and Code & 2026 & OA & qualifies or contradicts the claim (0.70) \\
Prompting for Security: A Cross-Model Evaluation of Code Generation in LLMs & 2025 & IEEE, OA, Scopus, S2 & qualifies or contradicts the claim (0.65) \\
Prompting techniques for secure code generation: a systematic investigation & 2025 & OA, arXiv, WoS, Scopus, S2 & qualifies or contradicts the claim (0.50) (new-technique) \\
Quality and Security Signals in AI-Generated Python Refactoring Pull Requests & 2026 & Scopus, OA, S2, arXiv & qualifies or contradicts the claim (0.65) \\
Replication Package - The Language of Security: How Prompt Syntax Shapes Secure Code Generation in Open LLMs & 2026 & OA & qualifies or contradicts the claim (0.65) \\
Rethinking the Evaluation of Secure Code Generation & 2025 & arXiv, WoS, OA, Scopus, S2 & qualifies or contradicts the claim (0.55) \\
SecDrift: Measuring Sector-Conditioned Security Drift in AI-Generated Code & 2026 & OA, arXiv, S2 & qualifies or contradicts the claim (0.85) \\
Secure Code Generation at Scale with Reflexion & 2025 & WoS, IEEE, OA, Scopus, S2, arXiv & qualifies or contradicts the claim (0.65) \\
Secure Coding with AI -- From Detection to Repair & 2025 & arXiv, OA, Scopus, WoS, S2 & qualifies or contradicts the claim (0.85) \\
Securing web apps with ChatGPT: A teaching plan for undergraduate students & 2026 & WoS, OA, Scopus & qualifies or contradicts the claim (0.70) \\
Security Impact of Coding Agent's Memory Files & 2026 & OA, arXiv, S2 & qualifies or contradicts the claim (0.55) (attack-on-model) \\
Security Implications of Large Language Model Code Assistants: A User Study & 2022 & OA, S2 & qualifies or contradicts the claim (0.95) \\
Security in the age of AI teammates: An empirical study of agentic pull requests on GitHub & 2027 & WoS, OA, Scopus, S2, arXiv & qualifies or contradicts the claim (0.80) \\
Security Knowledge Dilution in Large Language Models: How Irrelevant Context Degrades Critical Domain Expertise & 2025 & OA & qualifies or contradicts the claim (0.60) \\
Security Vulnerabilities in AI-Generated Code: A Large-Scale Analysis of Public GitHub Repositories & 2025 & WoS, OA, Scopus, S2, arXiv & qualifies or contradicts the claim (0.85) \\
Slopsquatting en código generado por modelos de lenguaje: análisis empírico de la alucinación de paquetes npm en el desarrollo web & 2026 & OA, S2 & qualifies or contradicts the claim (0.80) \\
The Effect of Security-Focused Prompts on Reported Findings: A Controlled Experiment on 400 AI-Generated Web Applications & 2026 & OA, arXiv, S2 & qualifies or contradicts the claim (0.80) \\
The Impact of AI-Assisted Development on Software Security: A Study of Gemini and Developer Experience & 2026 & OA, arXiv, S2 & qualifies or contradicts the claim (0.90) \\
The impact of GitHub Copilot on the Security of Code by Novice and Experienced Developers & 2025 & OA, S2 & qualifies or contradicts the claim (0.90) \\
The Range Shrinks, the Threat Remains: Re-evaluating LLM Package Hallucinations on the 2026 Frontier-Model Cohort & 2026 & OA, arXiv, S2 & qualifies or contradicts the claim (0.80) \\
Vulnerability Trends across Open-Source AI-Attributed Code & 2026 & OA, S2 & qualifies or contradicts the claim (0.75) \\
What is the Difference Between Me and You? Benchmarking the Quality Gap Between Human-Written and AI-Generated Code & 2026 & OA, S2, arXiv & qualifies or contradicts the claim (0.85) \\
When AI Writes Code: Investigating Security Issues in Agentic Software Changes & 2026 & OA, Scopus, S2 & qualifies or contradicts the claim (0.75) \\
\end{longtable}
\endgroup

